\section{Metric topology at all precisions}
\label{sec:metric-names}

The source article says that a metric defines a topology, gives compactness and
connectedness as examples of invariants, and compares a cube and a sphere.
The linked articles on metric spaces and compact spaces also discuss completeness,
uniform continuity, finite nets, and extrema. Here the points themselves can carry
infinitely many approximations. Each approximation, each accuracy request, and
each construction used to answer it has an arithmetic meaning.

\subsection{Points and balls}

Use a coded countable dense domain $D$ with a primitive-recursive (PR)
rational metric $d$ and its three metric laws. The examples below are rational
vectors with $d(u,v)=\max_i |u_i-v_i|$, including the one-dimensional case,
and rational vectors in a cube. Write $e_p=2^{-p}$. A \emph{name} is a fixed PR
recipe $x_p\in D$, with its numerical parameters suppressed, satisfying
\[
 d(x_p,x_q)\le e_p+e_q\qquad(p,q\ge0).
\]
For a rational $a\ge0$ abbreviate by $D(x,y)\le a$ the arithmetic sentence
\[
 \forall p,q\quad d(x_p,y_q)\le e_p+a+e_q .
\]
Equality of points means $D(x,y)\le0$. Apartness means that some pair
$p,q$ violates the last inequality with $a=0$. A rational point is named by
its constant recipe. Quantifiers over names, dense-domain operations, or
moduli throughout this section specify theorem schemes, with their displayed
numerical assumptions as antecedents. In particular no number is evaluated as
the code of an arbitrary PR function.

\begin{lemma}[Error removal, density, and the metric laws]\label{idx-metric32}
If a rational inequality $b\le c+C e_h$ holds for every sufficiently large
$h$, with a given rational $C\ge0$, then $b\le c$. The relations
$D(x,y)\le a$ are symmetric and satisfy the triangle inequality in the sense
of adding bounds. Equality is an equivalence relation. Each $x_p$ is within
$e_p$ of $x$.
\end{lemma}
\begin{proof}
Decide the rational comparison. If $b-c>0$, write $b-c=u/v$ with positive
integer numerator and denominator. Take $h$ at least the given threshold and
at least $v(\lceil C\rceil+1)+1$. Then $2^h>vC$, which contradicts
the proposed inequality. For the triangle inequality bridge through the
$h$th entry of the middle name. This adds $2e_h$ to the desired rational
bound at fixed outer indices. Remove that error by the first assertion.
Reflexivity is regularity of a name, and symmetry comes from $d$. For density,
$d(x_p,x_q)\le e_p+e_q$ is precisely the required estimate against a
constant name (one may remove the constant name's extra error).
\end{proof}

\begin{proposition}[A countable basis and separation]\label{idx-metric51}
Rational balls form a countable basis on the named completion. PR-indexed
unions of presented opens and finite intersections of presented opens are
presented opens. Apart points have disjoint rational neighborhoods, and
failure of apartness implies equality.
\end{proposition}
\begin{proof}
A certificate for $x\in B(c,r)$ is an index $p$ with
$d(x_p,c)+e_p<r$. Passing to a different name for $x$ preserves membership:
at a sufficiently large stage $h$, the excess error $2e_h$ is smaller
than the displayed rational margin. The same calculation works for closeness
to a constant point: if $D(x,c)\le b<r$, taking $2e_h<r-b$ gives a
membership certificate.

Suppose $x$ belongs to two rational balls, with certificates for each.
Let $m>0$ be the smaller of their two rational margins, and take $h$ with
$4e_h<m$. The ball $B(x_h,2e_h)$ contains $x$ and lies inside both balls:
for each original center $c$, the triangle inequality gives
$d(x_h,c)+2e_h\le d(x_p,c)+e_p+3e_h<r$.
A presented open is a PR-indexed list of balls with a bit enabling each
entry. The empty list and the list of all balls present the empty and whole
spaces. To intersect two lists, enumerate triples of two enabled entries
and a rational ball $B(c,s)$ satisfying $d(c,c_j)+s<r_j$ for both entries.
These rational tests are decidable, and the refinement just proved shows the
result is exact. Use pairing of indices for a PR double list to take its union;
iterate the binary construction for any fixed finite number of lists.
For a uniformly given finite input family one can instead code all original
indices in a single tuple and make all its containment tests in a bounded scan.

For apart points take a certificate with rational surplus
$m=d(x_p,y_q)-e_p-e_q>0$. Balls at $x_p,y_q$ of radii $e_p+m/3,e_q+m/3$
contain the respective points and are disjoint, since their radii sum to less
than the distance of their centers. If no certificate exists, each individual
rational test for a violation fails by decidability. These inequalities give
equality. The weaker negative statement that two names are unequal does not
supply the positive margin used in the separation construction.
\end{proof}

\subsection{Completion and extension from a dense domain}

\begin{theorem}[Completion with a Cauchy modulus]\label{idx-metric91}
Let $x_{i,p}$ be a PR double array of names and $M$ a PR modulus such that
$i,j\ge M(k)$ implies $D(x_i,x_j)\le e_k$. Then a PR name $z$ satisfies
$D(x_i,z)\le e_k$ for every $i\ge M(k)$ and is the unique limit up to
equality. Thus the represented Euclidean space, and the represented closed
cube, are complete for such sequences.
\end{theorem}
\begin{proof}
Replace $M(k)$ by $\max_{j\le k}M(j)$ for constructing the name. Call this
larger bound $M^*(k)$ and put $z_p=x_{M^*(p+1),p+1}$. For $p\le q$ the
outer indices both exceed or equal $M(p+1)$, so
$d(z_p,z_q)\le 2e_{p+1}+e_{q+1}\le e_p+e_q$. For the stated convergence
use any $i\ge M(k)$ and fixed $s,q$. Bridge from $x_{i,s}$ to $z_h$ and
then $z_q$, with $h\ge k$. The bound is
$e_s+e_k+e_{h+1}+e_h+e_q$. Remove the vanishing errors.
Two limits with moduli are within $2e_k$ for every $k$ by bridging through
an index at least both convergence thresholds, giving equality.
All entries remain in the dense cube if the input entries lie there. A
coordinatewise bounded name may also be clipped to the closed cube at each stage:
clipping is 1-Lipschitz and moves a stage by at most its certified error if
the named point belongs to the cube. Componentwise clipping fixes the named
point by error removal.
\end{proof}

For the next theorem the source and target may have different rational
metrics. A map on the dense source consists of a PR array $A(u,p)$ forming a
target name for each $u$, together with a nondecreasing PR modulus $\mu$ such
that
\[
 d(u,v)\le e_{\mu(k)}\quad\Longrightarrow\quad D(A(u),A(v))\le e_k .
\]
A nonmonotone supplied modulus can first be replaced by its running maximum.

\begin{theorem}[Uniform extension and open inverse images]\label{idx-metric124}
There is a unique uniformly continuous extension to names, up to equality.
Its modulus can be taken to be $\nu(k)=\mu(k)+1$. Its inverse images of
presented opens are presented opens. Identity and composition preserve this
class of maps, and compositions have the composed moduli.
\end{theorem}
\begin{proof}
Put $n(p)=1+\max(p,\mu(p+1))$ and
$F(x)_p=A(x_{n(p)},p+1)$. For $p\le q$ and equal names $x,y$ the arguments
$x_{n(p)},y_{n(q)}$ differ by at most $2e_{n(p)}\le e_{\mu(p+1)}$.
The two target samples differ by at most $2e_{p+1}+e_{q+1}\le e_p+e_q$.
Interchanging the roles handles the other order of indices. This proves both
regularity and independence of representatives. On constant names the
reindexing $p\mapsto p+1$ preserves equality by regularity.

If $D(x,y)\le e_{\mu(k)+1}$, then at all sufficiently large $h$ the two
arguments at stage $n(h)$ differ by at most
$e_{\mu(k)+1}+2e_{n(h)}\le e_{\mu(k)}$. Hence their output samples at
stage $h+1$ differ by at most $e_k+2e_{h+1}$. Bridge to arbitrary fixed
outer samples by regularity and remove the errors depending on $h$. For
uniqueness compare two candidate extensions at $x$ through a rational
approximant $x_h$ within both input thresholds for target error $e_k$.
They agree there, so their outputs at $x$ differ by at most $2e_k$ for all
$k$. For composition use the computed target name as input to the next
extension and compose its accuracy requests in reverse application order.

For an output ball $B(b,r)$ enumerate every source ball
$B(c,e_{\nu(k)})$ with $c$ rational for which some stage $p$ gives
\[
 d(A(c,p),b)+e_p+e_k<r.
\]
Both $k$ and $p$ occur in the enumeration index. Uniform continuity proves
the containment in the inverse image. Conversely, if $F(x)\in B(b,r)$
its certificate gives a positive margin $m$. Choose $3e_k<m$ and $c=x_h$
with $2e_h<e_{\nu(k)}$. Then $x$ belongs to the enumerated candidate
source ball and $D(F(c),F(x))\le e_k$. At a sufficiently large stage $p$,
$d(A(c,p),b)+e_p+e_k\le r-m+2e_k+2e_p<r$. This certifies the entry.
Take the double union for any presented output open.
\end{proof}

\subsection{Grids, covers, and extrema}

\begin{theorem}[Quantitative compactness of a cube]\label{idx-metric166}
For $Q=[0,1]^d$ in the maximum metric and every $k$, the explicit grid
\[
 G_k=\{0,1/2^{k+2},\ldots,1\}^d
\]
is an open $e_k$-net of $(2^{k+2}+1)^d$ points. A ball cover with a
supplied local selection at one uniform scale admits an explicit finite
subcover. Uniformly continuous images of the cube are totally bounded.
\end{theorem}
\begin{proof}
Round each coordinate of the clipped stage $x_{k+2}$ down to the mesh of
$G_k$ by integer division. The rational rounding costs at most $e_{k+2}$,
so the named point is within $e_{k+1}$ of the returned grid point. This
leaves a strict margin for radius $e_k$.
Here is the precise cover input. For every $g\in G_k$ a given PR selector
(or a supplied table) returns an enabled index $j(g)$ of a ball $B(c_j,r_j)$
and $d(g,c_j)+e_k\le r_j$. Retain this finite list of indices. For each
$x$, the preceding rounding supplies a $g$ with $D(x,g)\le e_{k+1}$.
The triangle inequality gives a strict membership certificate in the selected
ball, because $e_{k+1}<e_k$. An actual uniform Lebesgue-radius selector,
valid at all rational centers, supplies precisely such a table on the grid;
larger original open sets can receive balls with certified inclusion.
The proof uses the uniform scale supplied with the cover.

For an image net at radius $e_k$, apply the grid construction at input scale
$\nu(k+1)$ and map its centers through $F$. Every image point is within
$e_{k+1}$ of one of these finitely many image points. If a net in the target
dense domain is desired, replace each image center by its rational stage
$k+2$, leaving a strict margin. The resulting rational centers need not
themselves belong to the image.
\end{proof}

\begin{proposition}[Deciding rational box covers and computing a Lebesgue radius]\label{idx-metric198}
\label{prop:boxcover}
Given a finite list of rational open boxes $B_i=\prod_{j=1}^{d}(l_{ij},u_{ij})$,
there is a PR test for whether their restrictions cover $[0,1]^d$. Failure
returns a missed rational point. Success returns a rational $\lambda>0$
and a PR selection from each named point $x$ of one original box containing
the entire relative ball of radius $\lambda$ about $x$. This applies in
particular to finite lists of rational open balls in the maximum metric.
\end{proposition}
\begin{proof}
In dimension zero the cube and every empty product are singletons.
If the list is empty return the empty tuple as a missed point; otherwise
use the first box and $\lambda=1$. Now let $d\ge1$. Discard boxes with
any $l_{ij}\ge u_{ij}$, retaining the original indices of the others.
For each coordinate $j$, sort the distinct numbers in the list $E_j$
consisting of $0,1$ and all endpoints in that coordinate, including
endpoints outside $[0,1]$. Within $[0,1]$ take each breakpoint itself
and the midpoint of each pair of consecutive breakpoints as representatives.
Test every tuple of representatives for membership in some retained box,
using strict rational comparisons. If any tuple fails, return it.
Otherwise every rational point is covered: replace each of its coordinates
by the representative of the same singleton or open interval. All strict
comparisons with original endpoints are preserved, including comparisons
at a face or corner of the cube. If there are $b$ original boxes, at most
$(4b+3)^d$ tuples are tested. In particular an empty list fails.

On success let $4\delta$ be the smallest positive difference between
consecutive members of any $E_j$. There is such a difference since
$0,1\in E_j$. For a rational point $v$ in the cube, snap each coordinate
within $\delta$ of a breakpoint to that breakpoint, leaving other
coordinates alone, to form $v'$. The candidate is unique when it exists.
It lies in $[0,1]$ because every breakpoint outside that interval is at
least $4\delta$ away from the nearer endpoint $0$ or $1$. Scan the boxes
for one containing $v'$. For every snapped coordinate both boundaries of
this box are distinct breakpoints from the snapped value and hence at least
$4\delta$ away from it. At the original coordinate they are at least
$3\delta$ away. An unsnapped coordinate was already more than $\delta$
from every breakpoint. Thus on every coordinate $v$ has slack at least
$\delta$ inside this one box.

Return $\lambda=\delta/2$. For a named point $x$, take a clipped rational
sample $v$ at a fixed stage $h$ with $2^{-h}<\delta/4$ (for example
$h=1+\lceil4/\delta\rceil$). Apply the preceding selection to $v$.
Every named point of the relative ball about $x$ is within
$\delta/4+\delta/2<\delta$ of $v$ in the maximum metric and therefore
belongs to that box, with a positive coordinate margin. Sorting, finding
the minimum gap and scanning the boxes use only finite rational data.
\end{proof}

\begin{corollary}[Extremal values with near-optimizers]\label{idx-metric247}
For a real-valued map $F$ as above on a nonempty cube the supremum and
infimum exist as PR names. At precision $p$ one can return a rational point
$q_p$ with $F(x)\le F(q_p)+e_p$ for every $x$.
\end{corollary}
\begin{proof}
Use $G_{\nu(p+2)}$ and form the rational maximum $m_p$ of stage $p+2$
output samples on that grid, retaining any maximizing index $q_p$.
Every sample differs from its named value by at most $e_{p+2}$.
Uniform continuity and rounding thus give
$F(x)\le m_p+2e_{p+2}$ for every $x$, whereas
$m_p\le F(q_p)+e_{p+2}$. The two errors sum to less than $e_p$.
In particular $m_p\le m_q+e_{p+2}+2e_{q+2}$ and the reverse estimate
holds with $p,q$ interchanged. Thus $(m_p)$ is regular. The inequalities
show by error removal that its value bounds $F$ above and lies below every
upper bound. Apply the same argument to $-F$ for the infimum.
Here and below an inequality between two scalar names is read at stages,
with the sum of the two stage errors allowed. This construction provides
the extremal value and the requested near-optimizer at each precision;
it supplies no convergence modulus for the sequence of optimizing points.
\end{proof}

The line already has a PR open cover with no finite subcover:
$B(0,n+1)$ for $n\in\mathbb N$. A first rational sample and its error
bound give an index covering each named point. For any finite nonempty
selection with largest index $N$, the rational $N+2$ lies outside all
selected balls. The empty selection misses $0$ as well.

\subsection{Paths, a discrete test for connectedness, and intermediate values}

\begin{theorem}[Contractibility and a connectedness test]\label{idx-metric277}
The cube is path connected and contractible by maps with explicit moduli.
Any PR coloring of its rational points by two bits that is locally constant
at every named point is constant. This last assertion requires no uniform
modulus of local constancy.
\end{theorem}
\begin{proof}
For two named points $u,v$ use $t\mapsto (1-t)u+tv$ for $0\le t\le1$.
For rational $t$ compute the indicated convex combination on the same
stages of both names. Its error against another such stage is bounded by
the same convex combination of the two error bounds. Varying $t$ costs
at most $|t-t'|$ since each coordinate difference in the unit cube has
magnitude at most one. Extension gives the entire path. Fix an anchor $a$
and put $H(x,t)=(1-t)x+ta$. Changing $x$ costs at most its displacement
and changing $t$ costs at most its displacement. Thus $k+1$ is a modulus
for the product maximum metric at target precision $k$. The endpoints
are the identity and the constant map, and $H(a,t)=a$.

For the second assertion, local constancy means that at every named $z$
there exist a positive rational radius and a bit such that every rational
point within this radius of $z$ has that bit. If two rational endpoints
have different bits, bisect their segment and query its rational midpoint.
Keep that half whose endpoint bits still differ. All choices are comparisons
of rational-point bits. After $n$ steps the endpoints are at distance at
most $2^{-n}$, and subsequent segments are contained in this segment.
The sequence of left endpoints therefore forms a PR regular name $z$, and
both endpoints of the $n$th segment are within $2^{-n}$ of it. These
assertions follow by bounded induction on intermediate subdivision steps:
each retained segment is a subsegment, and its length halves. Choose $n$
with $2^{-n}$ smaller than the radius supplied at $z$. Both endpoints must
now have the same bit, a contradiction. Thus all rational points have one
bit. A rational approximant to any $z$ in its local constancy neighborhood
has that same bit. This proves the claimed test. It rules out clopen
partitions whose side on rational inputs has a PR decision procedure and
whose local membership is witnessed on all named points.
\end{proof}

\begin{theorem}[Intermediate values and a fixed point at prescribed accuracy]\label{idx-metric314}
Let $F:[0,1]\to\mathbb R$ be given by the dense-domain array and modulus,
and suppose $F(0)\le0\le F(1)$. A PR construction returns $t_p\in[0,1]
\cap\mathbb Q$ satisfying $|F(t_p)|\le e_p$. If the actual values of $F$
on rational inputs are given by a rational-valued PR function, one can
instead construct a single named point $t$ with $F(t)=0$.
Every uniformly presented self-map of the interval consequently has
approximate fixed points, and the rational-valued-on-rationals case has
an exact named fixed point.
\end{theorem}
\begin{proof}
Put $s=p+2$ and $N=2^{\mu(s)}$. At $i/N$ for $0\le i\le N$
take an output sample $v_i$ of precision $s$. If $v_0\ge0$ then
$-e_s\le F(0)\le0$. If $v_N<0$ then $0\le F(1)\le e_s$.
In the remaining case scan for the first $i$ with $v_i\ge0$.
Its predecessor has $v_{i-1}<0$, hence $F((i-1)/N)\le e_s$.
The mesh and the dense-domain modulus give $F(i/N)\le 2e_s$,
while $v_i\ge0$ gives $F(i/N)\ge-e_s$. Return the indicated sample
point. Subtracting a target value reduces other intermediate values to zero
provided its name and the shifted evaluation scheme are supplied.

For exact rational evaluation bisect $[0,1]$. Keep $[a,m]$ when the
midpoint value is nonnegative, and $[m,b]$ otherwise. Each step preserves
$F(a)\le0\le F(b)$ and halves the length. The left endpoints give a PR name
$t$, within $2^{-n}$ of both endpoints of the $n$th interval. For every
$k$ choose $n\ge\nu(k)$; extension continuity and the two endpoint
inequalities give $-e_k\le F(t)\le e_k$. Error removal proves equality.
In this bisection the comparison at a possible zero is a decidable
comparison of rationals. For the fixed-point assertions apply the
corresponding construction to $t-F(t)$ with the displayed endpoint signs.
For a dense-domain modulus $\mu$ of $F$, the difference has modulus
$k\mapsto\max(k+1,\mu(k+1))$.
\end{proof}

\subsection{A cube boundary, a round sphere, and the open interval}

\begin{theorem}[Radial homeomorphism, uniformly in the dimension]\label{idx-metric350}
For $d\ge1$ the closed cube $[-1,1]^d$ and the closed Euclidean unit ball
are homeomorphic. The restriction takes the cube boundary homeomorphically
onto the unit sphere. In maximum-metric coordinates the forward and
inverse maps have Lipschitz constants at most $d+2$ and $3d$ respectively.
In particular the boundary of a square is homeomorphic to a circle.
\end{theorem}
\begin{proof}
Put $a(x)=\max_i|x_i|$ and $s(x)=(\sum_i x_i^2)^{1/2}$. The maximum and
the sum are finite folds. A square root of a nonnegative rational $v$ is a
regular name: halve $[0,1+\lceil v\rceil]$ repeatedly, testing whether
the rational midpoint squared exceeds $v$, and take an endpoint once the
interval length is at most $e_p$. At most $p+1+\lceil v\rceil$ halvings
suffice. The nested intervals give a unique nonnegative value of square $v$,
since the error in the square is bounded by twice the initial upper bound
times the interval length. For a nonnegative named scalar the same procedure
applies to its nonnegative rational sample at stage $2p+4$, using the
inequality $|\sqrt{u}-\sqrt{v}|\le\sqrt{|u-v|}$ (square both sides after
ordering rational inputs). Computing the rational sample's square root to
precision $p+2$ leaves total error at most $2e_{p+2}$.
Multiplication on a bounded interval is PR on names: bounds from the zeroth
samples plus one control the errors in each factor. To divide by a positive
quantity with rational lower bound $b$, clip denominator samples below at
$b$; the reciprocal changes by at most $b^{-2}$ times the sample error.
Thus the arithmetic below has PR realizers with bounded precision requests.

For rational $x\ne0$ define
\[
 P(x)=a(x)\frac{x}{s(x)},\qquad T(x)=s(x)\frac{x}{a(x)} .
\]
Set both values to zero at the rational origin. Each nonzero rational input
has the explicit lower bound $a(x)>0$ for both denominators. To extend
across the origin and to arbitrary named inputs we prove uniform estimates.
For rational $x,y$ write $\delta=a(x-y)$. We have
$|a(x)-a(y)|\le\delta$ and $a(x)\le s(x)\le d\,a(x)$ by summing squares.
Furthermore
\[
 |s(x)^2-s(y)^2|
 \le\sum_i|x_i-y_i|(|x_i|+|y_i|)
 \le d\delta\,(s(x)+s(y)).
\]
When either rational vector is nonzero there is a rational positive lower
bound on the sum in parentheses, giving $|s(x)-s(y)|\le d\delta$ by
division. If both are zero the inequality is immediate.

For nonzero $x,y$ set $N_s(x)=x/s(x)$. Writing the difference with first
denominator $s(x)$ and using $|y_i|/s(y)\le1$ gives
\[
 a(N_s(x)-N_s(y))\le(\delta+|s(x)-s(y)|)/s(x).
\]
Split $P(x)-P(y)$ into
$a(x)(N_s(x)-N_s(y))+(a(x)-a(y))N_s(y)$. Since $a(x)\le s(x)$ this
costs at most $(d+2)\delta$. The identical subtraction with
$N_a(x)=x/a(x)$ gives $a(N_a(x)-N_a(y))\le2\delta/a(x)$.
Split $T(x)-T(y)$ using $s(x)\le d\,a(x)$ and the bound on
$|s(x)-s(y)|$ to get $3d\delta$. With one argument zero use
$a(P(x))\le a(x)$ and $a(T(x))=s(x)\le d\,a(x)$.
A Lipschitz constant $L\ge1$ gives the PR modulus $k+\lceil L\rceil$.
Consequently both maps extend to all named vectors.

At a vector with a witnessed positive lower bound for $a(x)$ the extended
formulas agree with the displayed quotients: nearby rational approximants
have half the given lower bound, so the reciprocal estimate proves equality
by error removal. For rational nonzero $x$ one has $s(P(x))=a(x)$ and
$a(P(x))=a(x)^2/s(x)>0$, and direct substitution gives $T(P(x))=x$.
Likewise $a(T(x))=s(x)>0$ and $P(T(x))=x$. At the rational origin both
identities hold by construction. Equality on the dense rational domain
extends to all names by the uniqueness argument above. The identities
$s(P(x))=a(x)$ and $a(T(x))=s(x)$ extend in the same way. They show
exactly what happens to the closed unit sets and their boundaries.
No zero test on a named vector was used.
\end{proof}

There is also a quantitative version of the article's distinction between
the circle and the line. The boundary of $[-1,1]^2$ has four rational
edge grids of mesh $e_{k+2}$. For a boundary name $x$, take its clipped
stage $k+3$. A coordinate of largest absolute rational value is within
$e_{k+3}$ of magnitude one because $a(x)=1$. Move that coordinate to its
signed endpoint and round the other coordinate down on the edge grid.
The resulting point is within $2e_{k+3}+e_{k+2}<e_k$ of $x$.
The radial map therefore supplies finite nets on the circle. Any real-valued
map on it with a uniform modulus has bounded image: take a finite net at
the input accuracy for target radius one, and bound the finitely many center
values by the largest absolute zeroth sample plus one. Adding one bounds
the whole image. Thus no such map can be a surjection onto the line.
The uniform-modulus hypothesis here belongs to the representation.


\begin{theorem}[Two charts for the sphere]\label{thm:stereographic}\label{idx-metric438}
For $n\ge1$, each sphere $S^n\subseteq\mathbb R^{n+1}$ has two explicit
Euclidean charts. Each chart identifies the complement of one pole,
with positive separation from that pole, with the whole of $\mathbb R^n$.
The charts cover the sphere with a finite decision at one fixed precision.
Every $S^n$ is path connected. Finite products of these spheres are
path connected and have product Euclidean charts; in particular the torus
$S^1\times S^1$ has two-dimensional charts.
\end{theorem}
\begin{proof}
Write $x=(v,z)$ with $\sum v_i^2+z^2=1$. For
$\sigma\in\{1,-1\}$ the domain $U_\sigma$ consists of points with
a witnessed positive lower bound on $1-\sigma z$. This describes
positive separation from the pole $(0,\sigma)$:
the squared Euclidean distance to that pole is $2(1-\sigma z)$, and
the maximum and Euclidean norms bound each other as proved above. A
smaller rational lower bound on the same expression persists on a small
neighborhood, so the domain is relatively open. Put
\[
 \phi_\sigma(v,z)=\frac{v}{1-\sigma z},\qquad
 \psi_\sigma(t)=
 \left(\frac{2t}{1+w},\frac{\sigma(w-1)}{1+w}\right),
 \qquad w=\sum_{i=1}^n t_i^2 .
\]
The identity $4w+(w-1)^2=(1+w)^2$ puts $\psi_\sigma(t)$ on the sphere.
Its denominator gap at the excluded pole is $2/(1+w)>0$, so substitution
gives $\phi_\sigma\psi_\sigma(t)=t$. In the other direction the sphere
equation gives
$\sum_i (v_i/(1-\sigma z))^2=(1+\sigma z)/(1-\sigma z)$.
Both coordinates of $\psi_\sigma\phi_\sigma(v,z)$ now simplify to
$(v,z)$.

For inputs to $\phi_\sigma$ with denominator at least $b>0$, the maximum
norm displacement of the output is at most $(b^{-1}+b^{-2})$ times the
input displacement. This follows by subtracting the quotients and using
$|v_i|\le1$. Given a gap at one point, use half of it on a neighborhood.
For inputs of maximum norm at most $R$, say $R\ge1$, the change in $w$
is at most $2nR$ times the input displacement. The first coordinates of
$\psi_\sigma$ then have Lipschitz bound $2+4nR^2$. For the last coordinate
write $(w-1)/(1+w)=1-2/(1+w)$ to get the bound $4nR$. These bounds
supply all the local moduli, and bounds for $R$ come from zeroth samples
plus their error. The chart changes are compositions of these explicit
maps on their overlap. The relative rational balls of the ambient space
supply a countable basis and separate apart points.

To cover a given $x$, inspect the last coordinate's rational sample at
stage $2$. If it is nonnegative take $\sigma=-1$, and otherwise take
$\sigma=1$. In either case $1-\sigma z\ge3/4$. The equatorial point
$e=((1,0,\ldots,0),0)$ lies in both charts with coordinate
$e_1=(1,0,\ldots,0)$. In the selected chart join $\phi_\sigma(x)$ to
$e_1$ by linear interpolation, then apply $\psi_\sigma$. A numerical
bound on the two endpoints bounds the whole segment and gives a uniform
modulus for this path by the preceding estimates. To join $x$ to $y$,
follow the path from $x$ to $e$ and then the reverse of the path from $y$
to $e$. On rational times the two formulas are evaluated on the two
half intervals. Each is Lipschitz after rescaling time, and at times on
opposite sides the two bounds add through the common value at $1/2$;
the larger Lipschitz constant works on the whole interval. Extend from
rational times as above. For a finite product take paths coordinatewise.
The tuple of chart choices has open product domain (use the minimum of
the finitely many local radii); concatenating coordinate lists identifies
the chart target with Euclidean space in the sum of the dimensions.
\end{proof}

\begin{proposition}[Completeness depends on the metric]\label{idx-metric502}
The open interval $(-1,1)$ and the line have an explicit homeomorphism
with local moduli. With their usual metrics the line is complete for named
Cauchy sequences with modulus, and the open interval is incomplete.
\end{proposition}
\begin{proof}
Take $f(t)=t/(1+|t|)$ and $g(u)=u/(1-|u|)$. The first has
$1-|f(t)|=1/(1+|t|)>0$, with a rational lower bound obtained from the
zeroth sample of $t$. A name in the open interval comes with a positive
rational bound on $1-|u|$. Thus both divisions are available and the
identities $g(f(t))=t$ and $f(g(u))=u$ follow by substitution.
Absolute value is 1-Lipschitz. On $|u|,|v|\le1-b$ with $0<b\le1$,
subtracting the two quotients bounds $|g(u)-g(v)|$ by
$(b^{-1}+b^{-2})|u-v|$. Given a local gap, use half of it on a
neighborhood to obtain a local modulus. On $|t|,|v|\le R$ the same
subtraction bounds $|f(t)-f(v)|$ by $(1+R)|t-v|$.
These bounds give open inverse image neighborhoods in both directions.

The rational sequence $1-e_n$ lies in the open interval and its terms
at indices $i,j\ge k$ are within $e_k$. Its unique limit on the line is
$1$, since the distance of the $n$th term from $1$ is $e_n$. A limit
in the open interval would yield the same line limit and simultaneously
a positive gap from $1$, a contradiction.
\end{proof}

\paragraph{Lean implementation order.}
Use rational metric domains with a decidable carrier code and coded vector
length; implement comparisons on pairs of integers by cross multiplication
with positive denominators. Introduce transparent predicates
\texttt{Reg x}, \texttt{DistLE x y r}, and \texttt{BallMem x c r} using
the formulas above. Prove the rational error-removal lemma first, then the
bridge triangle, evaluation against constants, and the diagonal completion
lemma. For maps use fields for the rational-stage evaluator, its regularity,
the modulus, and the displayed dense-domain bound. Define extension by the
index $n(p)$ above and prove its regularity before its modulus or its
inverse-image formula. For total boundedness implement the grid index in
base $2^{k+2}+1$ and prove the rounding estimate coordinatewise. For
intermediate values implement the finite scan and the exact-sign
interval recursion separately. Prove nestedness between two given stages
by induction on their gap, with the initial stage fixed and stage comparisons
as Boolean tests. For the radial equivalence first implement rational
square-root brackets and division with a lower bound, then the two estimates
on rational vectors and finally uniqueness of extension. A Lean layer
may quantify over function-valued parameters for reuse. Each iPRA instance
instead substitutes fixed PR terms, keeps regularity and modulus laws as
hypotheses at the numerical indices queried, and uses induction on finite
loops with decidable invariants.
